% =====================================================================
%  Apêndice B — Orçamento de energia (instrumento da Aula 08)
% =====================================================================
\chapter{Orçamento de energia}
\label{ap:energia}

O orçamento descreve o nó da AP2: sempre ligado, num carregador USB, com
a CPU a \SI{80}{\mega\hertz} e o BLE anunciando ou conectado o tempo
todo, num ciclo de \SI{300}{\second} (288 ciclos por dia). Os estados são
os da Seção~\ref{sec:firmware}. A inicialização acontece uma vez por
energização e não entra na conta; a nova tentativa tem o mesmo consumo de
\emph{mede}. Os LEDs de aviso da AP1 saíram. O nó não dorme.

\begin{table}[H]
\centering
\caption{Orçamento de energia do nó (alimentação USB)}
\label{tab:energia}
\scriptsize
\setlength{\tabcolsep}{3.5pt}
\begin{tabular}{@{}L{1.5cm}L{3.6cm}R{1.3cm}R{1.6cm}R{1.1cm}R{1.5cm}L{2.6cm}@{}}
\toprule
\textbf{Estado} & \textbf{O que está ligado} & \textbf{Cor\-rente (mA)}
& \textbf{Duração por ciclo (s)} & \textbf{Ciclos /dia}
& \textbf{Carga (mA$\cdot$h /dia)} & \textbf{Origem do número} \\
\midrule
aguarda & ESP32 a \SI{80}{\mega\hertz} com BLE (45), DHT11 em repouso (0,1), BH1750 desligado ($<$0,01), LED do módulo DHT11 (1,3)
  & 46,4 & 299,873 & 288 & 1.113,2 & \estimativa{} (README; \folhadedados) \\
mede & como \emph{aguarda}, com o DHT11 medindo (1,5) e o BH1750 convertendo (0,2)
  & 48,0 & 0,125 & 288 & 0,5 & \estimativa \\
envia & pico de transmissão BLE a \SI{0}{dBm} (130), por cerca de \SI{2}{\milli\second}
  & 131,4 & 0,002 & 288 & 0,02 & \folhadedados{} (pico); duração \estimativa \\
dorme & não existe: o nó não dorme & --- & --- & 0 & 0,0 & --- \\
\midrule
\multicolumn{5}{@{}r}{\textbf{Total diário orçado}} & \textbf{1.113,7} & \\
\multicolumn{5}{@{}r}{Corrente média medida (\CT{07})} & --- & não medida \\
\bottomrule
\end{tabular}
\fonte{Elaborada pela equipe. Corrente do rádio de \citeonline{espressif2023}.}
\end{table}

Origem de cada parcela:

\begin{itemize}
  \item \textbf{ESP32 com BLE, \SI{45}{\milli\ampere}:} topo da faixa de 35
    a \SI{45}{\milli\ampere} estimada no README para a CPU a
    \SI{80}{\mega\hertz} com o BLE ativo. É a linha que mudou de natureza
    na AP2: na AP1 era a CPU com o rádio desligado; agora é a CPU com o
    rádio ligado o tempo todo. Não foi medida: a equipe não tem medidor
    USB.
  \item \textbf{Rádio:} a transmissão em si quase não pesa. Um pacote de
    \SI{19}{\byte} ocupa o ar por frações de milissegundo; mesmo
    supondo \SI{2}{\milli\second} a \SI{130}{\milli\ampere}
    \cite{espressif2023}, são \SI{0.02}{\milli\ampere\hour} por dia. O que
    pesa é manter o rádio pronto para anunciar e para os eventos de
    conexão, e isso já está nos \SI{45}{\milli\ampere} de \emph{aguarda}.
    A linha do rádio deixou de ser a estimativa de WiFi da AP1 (pico de
    \SI{310}{\milli\ampere}): os tempos do enlace estão medidos na
    Tabela~\ref{tab:custo-radio}, e a corrente é a da folha de dados.
  \item \textbf{DHT11:} \SI{1.5}{\milli\ampere} medindo e
    \SI{0.1}{\milli\ampere} em repouso, ordem de grandeza da folha de dados
    \cite{aosongdht11}; LED do módulo, \SI{1.3}{\milli\ampere}, a mesma
    suposição da AP1.
  \item \textbf{BH1750:} cerca de \SI{0.12}{\milli\ampere} convertendo e
    \SI{1}{\micro\ampere} desligado \cite{rohmbh1750}; no modo de uma
    medida o chip se desliga sozinho depois de cada conversão.
\end{itemize}

\section*{Quatro linhas de decisão}
\phantomsection
\addcontentsline{toc}{section}{Quatro linhas de decisão}

\begin{description}[style=unboxed, leftmargin=0pt, font=\normalfont\bfseries]
  \item[1. Consumo diário total e consumo médio.]\hfill\\
    Orçado: \SI{1113.7}{\milli\ampere\hour} por dia, ou
    $1.113{,}7/24 \approx \SI{46.4}{\milli\ampere}$ de média. Na AP1 eram
    \SI{1577.2}{\milli\ampere\hour} (\SI{65.7}{\milli\ampere}): a saída dos
    três LEDs de aviso e do LED do fotorresistor tirou mais do que o BLE
    acrescentou, no papel. A comparação com uma medição fica para a AP3, porque a equipe não tem medidor USB (\CT{07}).
  \item[2. Fonte de alimentação escolhida, e por quê.]\hfill\\
    Carregador USB comum de \SI{5}{\volt} e \SI{1}{\ampere} ou mais, na
    tomada. O nó mede o ambiente de plantas dentro de casa, perto de uma
    tomada. A média de \SI{46}{\milli\ampere} é cerca de 9\,\% de uma porta
    USB 2.0 e o pico de \SI{130}{\milli\ampere}, cerca de 26\,\%. Um
    carregador rápido foi descartado porque religava a saída com carga
    baixa (Capítulo~\ref{cap:testes}, \enquote{O que falhou}); pelo mesmo
    motivo, muitos \emph{power banks} desligam sozinhos abaixo de 50 a
    \SI{100}{\milli\ampere}.
  \item[3. Autonomia estimada, e a conta que levou até ela.]\hfill\\
    Na tomada, ilimitada. Como referência, uma 18650 de
    \SI{2600}{\milli\ampere\hour} com 80\,\% de profundidade de descarga dá
    \SI{2080}{\milli\ampere\hour} úteis, e $2.080 \div 1.113{,}7 \approx
    1{,}9$ dia. O regulador da placa é linear, então a corrente da bateria
    é praticamente a da carga.
  \item[4. O que acontece quando a energia acaba --- ou enquanto o nó dorme.]\hfill\\
    O nó apaga por completo e perde a luz acumulada. Não há atuador nem
    estado a proteger. Diferente da AP1, a falta de energia agora é
    percebida por quem está longe do nó: o gravador anota
    \texttt{desconectado} em poucos segundos (\CT{18}), e, quando a energia
    volta, o id do reinício novo mostra no conjunto de dados que houve
    queda. O nó não dorme.
\end{description}

\section*{Se a conta atravessar tensões diferentes}
\phantomsection
\addcontentsline{toc}{section}{Se a conta atravessar tensões diferentes}

Não é o caso: o ESP32 e os dois módulos estão no barramento de
\SI{3.3}{\volt}, alimentado pelo regulador linear da placa a partir dos
\SI{5}{\volt} da USB. Num regulador linear a corrente de entrada é a de
saída mais a corrente própria do regulador, então os
\si{\milli\ampere\hour} da tabela valem também do lado dos
\SI{5}{\volt}.

\section*{Se o projeto for alimentado por painel}
\phantomsection
\addcontentsline{toc}{section}{Se o projeto for alimentado por painel}

Não se aplica: o nó fica dentro de casa, na tomada.

\section*{Quem está dimensionando este sistema}
\phantomsection
\addcontentsline{toc}{section}{Quem está dimensionando este sistema}

O ESP32 acordado com o rádio responde por cerca de 97\,\% do consumo
orçado ($45 \times 24 = \SI{1080}{\milli\ampere\hour}$ de
\SI{1113.7}{\milli\ampere\hour}); o LED do módulo DHT11 por cerca de 3\,\%;
os dois sensores, juntos, por menos de 0,1\,\%. A transmissão dos pacotes
não aparece na conta. A decisão que dimensiona o sistema é a \DP{11}, com
o rádio ligado o tempo todo à espera do central; uma versão a bateria
teria de dormir entre as medidas e anunciar só por alguns segundos.
